\section{Impact of Tool Augmentation on Single-Agent Baselines}
\label{sec:appendix_tool_ablation}

\subsection{Motivation and Setup}
A potential concern regarding the performance gains of DynaDebate is whether the improvements stem primarily from the integration of external tools rather than the proposed multi-agent debate mechanism itself. To rigorously disentangle the contribution of the debate framework from the benefits of tool utilization, we conducted additional comparative experiments using tool-augmented single-agent baselines.

\subsection{Implementation Details}
We established a strong baseline, denoted as Tool-Augmented CoT, which equips the standard Chain-of-Thought (CoT) method with external utilities. Crucially, to ensure a fair comparison, the tool availability for the baseline mirrors the configuration of the Verification Agent in our main experiments:

\begin{itemize}
    \item For MATH500: The baseline is equipped with a Python Code Interpreter (CoT + Code). It is prompted to write and execute Python code to perform calculations or verify logic steps.
    \item For MMLU: The baseline is equipped with a Search Engine (CoT + Search). It is allowed to query external information to retrieve factual knowledge before generating the final answer.
\end{itemize}


\subsection{Results}
The comparative results between the standard single-agent baseline (Standard CoT), the tool-augmented baseline (Tool-Augmented CoT), and our proposed DynaDebate framework are presented in Table~\ref{tab:tool_ablation}.

\begin{table}[h]
\centering
\small
\begin{tabular}{lcc}
\toprule
\textbf{Method} & \textbf{MATH500} & \textbf{MMLU} \\
\midrule
\multicolumn{3}{l}{\textit{\textbf{GPT-4o-mini}}} \\
CoT & 70.20 & 82.67 \\
\textit{Tool-Augmented CoT} & 70.60 & 83.00 \\
DynaDebate (Ours) & \textbf{73.60} & \textbf{83.33} \\
\midrule
\multicolumn{3}{l}{\textit{\textbf{Qwen3-8B}}} \\
Standard CoT & 78.00 & 78.67 \\
\textit{Tool-Augmented CoT} & 77.20 & 81.33 \\
DynaDebate (Ours) & \textbf{83.80} & \textbf{83.00} \\
\bottomrule
\end{tabular}
\caption{Performance comparison with tool-augmented single-agent baselines across distinct backbones.}
\label{tab:tool_ablation}
\end{table}

\noindent\textbf{Analysis of Results}

The results in Table~\ref{tab:tool_ablation}, considered alongside our main ablation studies (Table~\ref{tab:ablation}), provide a nuanced understanding of the relationship between tool availability and reasoning performance:

\begin{enumerate}
    \item \textbf{Tools are not a panacea; improper usage can hinder performance.}
    A critical observation from the Qwen3-8B results on MATH500 is that naively equipping a single agent with a code interpreter actually \textit{decreased} accuracy from 78.00\% to 77.20\%. This indicates that without a structured framework to verify and critique code execution, agents are susceptible to "tool misuse" (e.g., generating erroneous code or misinterpreting outputs), which introduces noise rather than clarity. DynaDebate avoids this by using tools strictly under triggered conditions within a peer-review setting, achieving 83.80\%.

    \item \textbf{Synergy between Search and Debate exceeds Search alone.}
    On the knowledge-intensive MMLU benchmark, we acknowledge that the Search Engine provides a tangible benefit to the single-agent baseline (e.g., Qwen3-8B improved from 78.67\% to 81.33\%). However, DynaDebate still achieves a superior performance (83.00\%). This demonstrates that our performance gains are not solely derived from accessing external information, but from the \textbf{DynaDebate framework's ability to effectively synthesize, debate, and verify} that information. The combination of search retrieval with multi-agent critique yields better outcomes than retrieval alone.

    \item \textbf{Architectural superiority persists without tools.}
    Finally, it is crucial to reference our ablation study in Section~\ref{sec:ablation} (Table~\ref{tab:ablation}). Even when the Verification module (and thus all tool access) is removed, the core DynaDebate framework (consisting only of Dynamic Path Generation and Debate) still outperforms the standard CoT baseline on benchmarks like MATH500 (81.80\% vs. 78.00\% for Qwen3-8B). This confirms that while tools act as a powerful amplifier, the fundamental reasoning improvements are driven by the proposed multi-agent mechanisms.
\end{enumerate}
